% =======================================================================
\chapter{Conclusión}
\label{cap:conclusion}
% =======================================================================

\section{El trabajo realizado}

El trabajo partió de una brecha concreta: la literatura sobre Búsqueda Tabú
compara familias de metaheurísticas o propone variantes especializadas, pero
rara vez aísla la contribución de los mecanismos internos de TS sobre problemas
heterogéneos. Para cerrarla se construyó un diseño factorial completo —cinco
problemas por tres niveles de TS, quince módulos, 360 experimentos— y se
publicó íntegramente como repositorio reproducible.

El resultado son tres productos complementarios, distintos en formato pero
derivados del mismo código:

\begin{enumerate}
  \item \textbf{El repositorio} (Capítulo~\ref{cap:repositorio}): quince
    implementaciones de Búsqueda Tabú, la infraestructura de orquestación,
    los modelos exactos de verificación y los resultados crudos de las 360
    corridas.
  \item \textbf{El artículo} (Capítulo~\ref{cap:paper}): el análisis formal
    en formato IEEE de conferencia.
  \item \textbf{El afiche} (Capítulo~\ref{cap:afiche}): la comunicación del
    trabajo en EXPOTIC UNAB 2026.
\end{enumerate}

\section{Los hallazgos}

\subsection{La progresión funciona, pero no de forma uniforme}

El GAP global promedio cae de \SI{6.50}{\percent} en TS-Básico a
\SI{4.00}{\percent} en TS-Intermedio y \SI{3.40}{\percent} en TS-Avanzado:
una reducción del \SI{48}{\percent} relativo entre los extremos. La
progresión es monótona en el agregado, lo que confirma la hipótesis de
partida.

El promedio, sin embargo, es engañoso. Desagregado por problema, casi toda la
ganancia proviene de CVRP (\SI{19.22}{\percent} $\rightarrow$
\SI{6.38}{\percent}, un \SI{67}{\percent} de reducción relativa), mientras
que en Flowshop los tres niveles quedan en torno al \SI{1.7}{\percent} y
la diferencia entre ellos carece de significado práctico.

\subsection{Los mecanismos rinden donde el espacio de búsqueda lo permite}

La heterogeneidad tiene una explicación estructural. Los mecanismos de
memoria de largo plazo —frecuencia, penalización adaptativa, diversificación
informada— sirven para escapar de óptimos locales; su valor depende de la existencia de óptimos locales de donde escapar. En Flowshop, la vecindad
de \emph{swap} adyacente limita lo que cualquier nivel de TS puede lograr
sobre la solución de la heurística NEH: los tres niveles terminan
prácticamente en el mismo GAP. En Knapsack, el \emph{greedy} por ratio
deja la solución tan cerca del óptimo que el margen de mejora es casi nulo.

CVRP~\cite{laporte1992vehicle} y Job Shop, en cambio, tienen espacios amplios y
accidentados, y es allí donde la memoria de largo plazo justifica su costo.

\subsection{El primer salto es el más rentable}

TS-Intermedio captura 2.50 puntos porcentuales de la mejora total;
TS-Avanzado añade solo 0.60 puntos adicionales. Considerando que
TS-Avanzado requiere aproximadamente $3.9\times$ el tiempo de cómputo de
TS-Básico, la relación costo-beneficio favorece claramente al nivel
intermedio en escenarios con restricción de tiempo. Este es, probablemente,
el hallazgo de mayor valor práctico del trabajo.

\subsection{La perturbación estructural importa}

El reinicio \emph{double-bridge} (4-opt)~\cite{lin1973effective} de TS-Avanzado
produjo la mayor mejora individual observada: \SI{2.21}{\percent} de reducción
en \texttt{kroA100}. Perturbar de forma informada la mejor solución conocida
resultó más eficaz que ampliar la lista tabú o alargar la búsqueda.

\section{Sobre el método}

Dos aspectos metodológicos merecen destacarse por encima de los resultados
numéricos.

El primero: el GAP se calcula contra \textbf{valores de referencia
verificables}. Los BKS de Knapsack se comprueban con programación dinámica,
y las instancias de Flowshop se regeneran desde las semillas publicadas por
Taillard, con sus \emph{makespan} óptimos como BKS. Con referencias de este
tipo ningún GAP puede ser negativo, y así ocurre en los 360 experimentos.

El segundo: la \textbf{validación sobre cinco semillas} resultó
indispensable. Con una sola semilla, varias de las diferencias entre niveles
habrían sido indistinguibles del ruido. La baja desviación estándar
observada es lo que permite afirmar que la progresión es real.

\section{Limitaciones}

El trabajo tiene fronteras claras, que conviene enunciar sin atenuantes:

\begin{itemize}
  \item \textbf{Criterio de parada por iteraciones.} Igualar iteraciones y no
    tiempo de CPU hace que las comparaciones de costo sean aproximadas. Una
    comparación estricta requeriría presupuesto temporal fijo.
  \item \textbf{Escala moderada.} Las 24 instancias van de 50 a 150
    elementos. Nada de lo aquí observado permite afirmar cómo escalan los
    mecanismos a instancias de orden de magnitud superior.
  \item \textbf{Un solo diseño por nivel.} Cada nivel es una configuración
    concreta, no una familia. Un TS-Básico mejor ajustado podría acercarse a
    TS-Intermedio.
  \item \textbf{Penalización simple.} La función de penalización por
    frecuencia emplea normalización $\phi/t$; esquemas más elaborados —por
    ejemplo, normalización por tamaño de vecindario— podrían mejorar la
    estabilidad.
  \item \textbf{Hardware limitado.} Los recursos disponibles condicionaron el
    número de iteraciones en las instancias mayores de CVRP y Job Shop.
\end{itemize}

\section{Trabajo futuro}

\begin{itemize}
  \item Sustituir el criterio de parada por presupuesto de tiempo de CPU fijo
    y repetir la comparación.
  \item Incorporar modelos exactos MILP para las instancias menores y
    contrastar contra cotas inferiores demostradas.
  \item Implementar un cuarto nivel con Búsqueda Tabú poblacional
    (trayectorias paralelas).
  \item Añadir instancias de mayor escala —TSP con más de 500 ciudades, CVRP
    con más de 500 clientes— y medir el escalado del GAP.
  \item Extender el repositorio a otros problemas: coloración de grafos,
    mochila multidimensional.
  \item Enriquecer la vecindad de Job Shop con movimientos sobre el grafo de
    disyunción~\cite{balas1998one}, dado que es el problema donde los GAP se
    mantienen más altos.
  \item Automatizar la generación de figuras desde los cuadernos Jupyter e
    integrarla a la compilación mediante \texttt{latexmk}.
\end{itemize}

\section{Cierre}

La arquitectura progresiva de Búsqueda Tabú validada en este trabajo permite
observar el impacto incremental de cada mecanismo metaheurístico en lugar de
tratarlos como una caja única. El hallazgo central no está en la victoria de la configuración avanzada —eso era esperable— sino \emph{cuánto} gana, \emph{dónde} gana
y \emph{a qué costo}: la mayor parte de la mejora se obtiene en el primer
salto de complejidad, se concentra en los problemas de espacio de búsqueda
amplio, y el nivel más sofisticado cobra casi cuatro veces el tiempo por cerca
de la mitad del GAP.

El código publicado permite que estas afirmaciones se verifiquen, se
cuestionen y se extiendan; esa es, en definitiva, la razón por la cual el repositorio —y no el artículo— es el producto principal del trabajo.
